\documentclass[10pt,a4paper]{amsart}
\usepackage[margin=19mm]{geometry}
\usepackage{amsmath,amssymb,mathtools}
\usepackage[scr=rsfs,cal=euler]{mathalfa}
\usepackage{xfrac}
\usepackage[unicode,hidelinks,bookmarks=false]{hyperref}
\newcommand{\ie}{i.e.\ }
\newcommand{\cf}{cf.\ }
\newcommand{\resp}{respectively\ }
\newcommand{\Subsection}{\subsection}
\providecommand{\nobreakdash}{\nobreak\hbox{-}}
\setlength{\emergencystretch}{2em}
\multlinegap0pt
\usepackage{amssymb}
\usepackage{mathtools}
\usepackage{bm}
\usepackage{xcolor}
\newtheorem{lemma}{Lemma}
\newtheorem{proposition}{Proposition}
\newcommand{\Fil}{\mathrm{Fil}}
\newcommand{\bR}{\mathbb{R}}
\newcommand{\fa}{\mathfrak{a}}
\newcommand{\fb}{\mathfrak{b}}
\newcommand{\fc}{\mathfrak{c}}
\title{On the geometry of spaces of filtrations on local rings}
\author{Lu Qi}
\date{Source: arXiv:2409.01705v2 (CC BY 4.0)}
\begin{document}
\maketitle

\noindent\emph{Source and licence.} On the geometry of spaces of filtrations on local rings. Version arXiv:2409.01705v2 (CC BY 4.0), distributed by arXiv under the Creative Commons Attribution 4.0 International licence, http://creativecommons.org/licenses/by/4.0/, as recorded in the arXiv licence metadata for this exact version and shown on the abstract page https://arxiv.org/abs/2409.01705v2. Full licence notice: \texttt{licenses/arxiv-2409.01705-CC-BY-4.0.txt}.\par\smallskip

\noindent\textit{Editorial scope.} Counted unit: the article's proposition minimizing (source0.tex lines 1555--1562). The article's complete original proof of it is delivered with it (source0.tex lines 1564--1606, one \texttt{proof} environment of 43 lines). No mathematics is written, repaired or re-derived: the statement and the proof are the article's own lines, in the article's own notation, and no retained line is substituted or re-flowed.  Retained uncounted as the source-local context the proof needs: lines 1526--1528.  Nothing was omitted from any retained span, so the package carries no omission record.  No retained line contains a citation, so the wrapper carries no bibliography and no bibliography entry.  No retained line contains a figure environment, an image-inclusion command or drawing code, so no image is delivered and the package declares no asset dependency.  Editorial formatting only, in addition: an AMS \texttt{amsart} preamble that restates the article's own declarations and macros used by the retained lines, namely the environments \texttt{lemma}, \texttt{proposition} on separate counters, together with the standard packages those lines need.  The retained environments print in this document as Lemma 1, Proposition 1 in order of appearance, whereas the article prints the same items with the numbering of its own full document.  The complete native source of the article as distributed in the arXiv e-print for this exact version is bundled unchanged as \texttt{sources/164667/source0.tex}.
\begin{lemma}\label{lem:subinterval}
    Let $\fa_{\bullet,i}\in\Fil$ for $i=0,1$, and let $\{\fa_{\bullet,t}\}_{t\in[0,1]}$ be the geodesic between them. For $0\le t<t'\le 1$, we have $\fa_{\bullet,0}\cap\fa_{\bullet,t'}\subset\fa_{\bullet,t}$. 
\end{lemma}

\begin{proposition}\label{prop:distance minimizing}
    Let $\fa_{\bullet,0},\fa_{\bullet,1}\in\Fil$ and let $\{\fa_{\bullet,t}\}_{t\in[0,1]}$ be the geodesic between them. Then for any $0\le t<t'\le 1$, we have
    \begin{equation}\label{eqn:distance minimizing}
        d_1(a_{\bullet,0},\fa_{\bullet,t'})=
        d_1(\fa_{\bullet,0},\fa_{\bullet,t})+
        d_1(\fa_{\bullet,t},\fa_{\bullet,t'}).
    \end{equation}
\end{proposition}

\begin{proof}
    For simplicity, denote $\fb_\bullet:=\fa_{\bullet,0}\cap\fa_{\bullet,t}$, $\fb'_\bullet:=\fa_{\bullet,0}\cap\fa_{\bullet,t'}$ and $\fc_\bullet:=\fa_{\bullet,t}\cap\fa_{\bullet,t'}$. By Lemma \ref{lem:subinterval} we know that $\fb'_\bullet\subset\fb_\bullet\subset\fa_{\bullet,0}$ and $\fb'_\bullet\subset\fc_\bullet\subset\fa_{\bullet,t'}$. We claim that for any $\lambda\in\bR_{>0}$, 
    \begin{equation}\label{eqn:parallelogram}
        \fa_{\lambda,t}/\fb_\lambda\cong \fc_\lambda/\fb'_\lambda.
    \end{equation}
	
    Granting the claim for now, we know that $d_1(\fb_\bullet,\fa_{\bullet,t})=d_1(\fb'_\bullet,\fc_\bullet)$. So we get
    \[
        d_1(\fb_\bullet,\fa_{\bullet,t})+
        d_1(\fa_{\bullet,t},\fc_\bullet)=
        d_1(\fb'_\bullet,\fc_\bullet)+
        d_1(\fa_{\bullet,t},\fc_\bullet)
        =d_1(\fa_{\bullet,t},\fb'_\bullet)
        =d_1(\fa_{\bullet,t},\fb_\bullet)+
        d_1(\fb_\bullet,\fb'_\bullet),
    \]
    which implies $d_1(\fb_\bullet,\fb'_\bullet)=d_1(\fa_{\bullet,t},\fc_\bullet)$. Now we compute
    \begin{align*}
        d_1(a_{\bullet,0},\fa_{\bullet,t'})=
        &d_1(\fa_{\bullet,0},\fb'_\bullet)+
        d_1(\fb'_\bullet,\fa_{\bullet,t'})\\
        =&d_1(\fa_{\bullet,0},\fb_\bullet)+d_1(\fb_\bullet,\fb'_\bullet)+d_1(\fb'_\bullet,\fc_\bullet)+d_1(\fc_\bullet,\fa_{\bullet,t'})\\
        =&d_1(\fa_{\bullet,0},\fb_\bullet)+d_1(\fa_{\bullet,t},\fc_\bullet)+d_1(\fb_\bullet,\fa_{\bullet,t})+d_1(\fc_\bullet,\fa_{\bullet,t'})\\
        =&d_1(\fa_{\bullet,0},\fa_{\bullet,t})+d_1(\fa_{\bullet,t},\fa_{\bullet,t'})
        \end{align*}
	
    It remains to prove the claim \eqref{eqn:parallelogram}. Indeed, by definition, any $f\in\fa_{\lambda,t}$ can be written as
    \[
        f=\sum_{(1-t)\mu+t\nu=\lambda}f_{\mu,\nu},
    \]
    where $f_{\mu,\nu}\in\fa_{\mu,0}\cap\fa_{\nu,1}$. If $\mu\ge \lambda$, then $f_{\mu,\nu}\in\fa_{\mu,0}\subset \fa_{\lambda,0}$. If $\mu\le\lambda$, then $\nu\ge \lambda$, and hence $f_{\mu,\nu}\in\fa_{\nu,1}\subset \fa_{\lambda,1}$. Thus we have 
    \[
        f=\sum_{\mu\le\lambda}f_{\mu,\nu}+
        \sum_{\mu\ge\lambda}f_{\mu,\nu}\in
        \fa_{\lambda,t}\cap\fa_{\lambda,1}+\fa_{\lambda,0},
    \]
    and hence $\fc_\lambda\subset\fa_{\lambda,t}\subset\fa_{\lambda,t}\cap\fa_{\lambda,1}+\fa_{\lambda,0}\subset\fc_\lambda+\fa_{\lambda,0}$, where the last inclusion follows from Lemma \ref{lem:subinterval}. So we get
    \[
        \fa_{\lambda,t}/\fb_\lambda\cong
        (\fa_{\lambda,t}+\fa_{\lambda,0})/\fa_{\lambda,0}=(\fc_\lambda+\fa_{\lambda,0})/\fa_{\lambda,0}\cong\fc_\lambda/(\fc_\lambda\cap\fa_{\lambda,0})=\fc_\lambda/\fb'_\lambda,
    \]
    where the last equality follows again from Lemma \ref{lem:subinterval}. This proves \eqref{eqn:parallelogram}, and finishes the proof of the proposition.
\end{proof}

\end{document}
